\documentclass[12pt,a4paper]{article}
\usepackage[utf8]{inputenc}
\usepackage[indonesian]{babel}
\usepackage{amsmath}
\usepackage{amsfonts}
\usepackage{amssymb}
\usepackage{geometry}
\usepackage{hyperref}

\geometry{margin=2.5cm}

\title{Metode Faktorisasi Polinomial Iteratif Berbasis Koefisien \\ (Metode Lin Derajat 4 dan 5)}
\author{Ahmad Fakhrul Bawani}
\date{}

\begin{document}

\maketitle

\section{Pendahuluan}
Metode Lin merupakan metode numerik iteratif untuk memfaktorkan polinomial derajat tinggi dengan cara mencocokkan koefisien hasil kali faktor-faktornya terhadap polinomial asli. Berbeda dengan metode rumpun Newton-Raphson, metode ini tidak membutuhkan perhitungan turunan (matriks Jacobian) maupun tebakan awal yang spesifik, melainkan memanfaatkan skema iterasi titik tetap (\textit{fixed-point iteration}) dengan nilai awal berupa nol.

\section{Formulasi untuk Polinomial Derajat 4}
Misalkan terdapat polinomial derajat 4 dengan bentuk umum:
\begin{equation}
P(x) = x^4 + A_3x^3 + A_2x^2 + A_1x + A_0
\end{equation}
Polinomial ini akan difaktorkan menjadi dua buah bagian kuadratik:
\begin{equation}
P(x) = (x^2 + b_1x + b_0)(x^2 + a_1x + a_0)
\end{equation}

Dengan mengalikan kedua faktor tersebut dan menyamakan koefisiennya dengan persamaan asli, diperoleh sistem persamaan non-linear sebagai berikut:
\begin{align}
A_3 &= a_1 + b_1 \\
A_2 &= a_0 + a_1b_1 + b_0 \\
A_1 &= a_1b_0 + a_0b_1 \\
A_0 &= a_0b_0
\end{align}

Sistem di atas ditransformasikan menjadi skema iterasi balik dengan mengisolasi masing-masing variabel:
\begin{align}
b_0 &= \frac{A_0}{a_0} \\
b_1 &= \frac{A_1 - a_1b_0}{a_0} \\
a_1 &= A_3 - b_1 \\
a_0 &= A_2 - b_0 - a_1b_1
\end{align}

\section{Formulasi untuk Polinomial Derajat 5}
Untuk polinomial derajat 5 dengan bentuk umum:
\begin{equation}
P(x) = x^5 + A_4x^4 + A_3x^3 + A_2x^2 + A_1x + A_0
\end{equation}
Faktorisasi dilakukan dengan memecahnya menjadi satu faktor linear dan dua faktor kuadratik:
\begin{equation}
P(x) = (x + a_0)(x^2 + b_1x + b_0)(x^2 + c_1x + c_0)
\end{equation}

Melalui ekspansi aljabar dan penyamaan koefisien secara runtut dari pangkat tertinggi hingga terendah, kita dapat menyusun urutan rumus pembaruan variabel untuk setiap siklus iterasi:
\begin{align}
c_1 &= A_4 - a_0 - b_1 \\
c_0 &= A_3 - a_0A_4 + a_0^2 - b_0 - c_1b_1 \\
b_1 &= \frac{A_2 - a_0c_0 - b_0c_1}{c_0} \\
b_0 &= \frac{A_1 - a_0b_1c_0 - b_0c_1}{c_0} \\
a_0 &= \frac{A_0}{b_0c_0}
\end{align}

\section{Algoritma Iterasi}
Prosedur komputasi dilakukan tanpa memerlukan input tebakan numerik dari luar, dengan langkah sebagai berikut:
\begin{itemize}
  \item \textbf{Langkah 1:} Inisialisasi seluruh variabel pembantu ($a_i, b_i, c_i$) dengan nilai awal \textit{starter} sama dengan $0$.
  \item \textbf{Langkah 2:} Untuk derajat 4, hitung nilai $b_0, b_1, a_1,$ dan $a_0$ berturut-turut menggunakan persamaan pada Seksi 2. Untuk derajat 5, hitung nilai $c_1, c_0, b_1, b_0,$ dan $a_0$ berturut-turut menggunakan persamaan pada Seksi 3.
  \item \textbf{Langkah 3:} Gunakan nilai yang baru saja diperoleh untuk dimasukkan kembali ke langkah perhitungan berikutnya.
  \item \textbf{Langkah 4:} Ulangi proses siklus tersebut hingga selisih nilai antar-iterasi (\textit{error}) bernilai sangat kecil atau konvergen.
\end{itemize}

\section{Analisis Metode}
\subsection{Kelebihan}
\begin{itemize}
  \item Struktur algoritma sangat sederhana dan linier, sehingga mudah diimplementasikan ke dalam kode program seperti C++ atau Python hanya dengan \textit{looping} biasa.
  \item Tidak memerlukan tebakan nilai awal yang rumit; nilai awal $0$ sudah cukup untuk memulai proses iterasi efek domino koefisien.
\end{itemize}

\subsection{Kekurangan}
\begin{itemize}
  \item Kecepatan konvergensi bertipe linier, cenderung lebih lambat dibandingkan metode Bairstow (Newton-Raphson) yang bertipe kuadratik.
  \item Rentan terjadi kegagalan operasi jika pada tengah-tengah iterasi nilai variabel pembagi ($a_0$ pada derajat 4 atau $c_0$ pada derajat 5) bernilai mendekati nol.
\end{itemize}

\end{document}